\documentclass[11pt]{article}

\usepackage[T1]{fontenc}
\usepackage{lmodern}
\usepackage[margin=1in]{geometry}
\usepackage{amsmath,amssymb,amsthm}
\usepackage{booktabs}
\usepackage{xcolor}
\usepackage{listings}
\usepackage[numbers,sort&compress]{natbib}
\usepackage[hidelinks]{hyperref}

% --- table marks -----------------------------------------------------------
\newcommand{\cmark}{\textcolor{black}{\checkmark}}
\newcommand{\xmark}{\textcolor{gray}{--}}
\newcommand{\tni}{\textsc{tni}}
\newcommand{\lt}{$\lambda_{\text{time}}$}

% --- theorem environments --------------------------------------------------
\theoremstyle{plain}
\newtheorem{theorem}{Theorem}
\theoremstyle{definition}
\newtheorem{definition}{Definition}

% --- listings --------------------------------------------------------------
\definecolor{kw}{rgb}{0.1,0.1,0.6}
\definecolor{cmt}{rgb}{0.3,0.5,0.3}
\lstdefinestyle{py}{
  basicstyle=\ttfamily\small, keywordstyle=\color{kw}, commentstyle=\color{cmt},
  language=Python, showstringspaces=false, breaklines=true, frame=single,
  framesep=4pt, columns=fullflexible}
\lstdefinestyle{plain}{basicstyle=\ttfamily\footnotesize, frame=single,
  framesep=4pt, breaklines=true, columns=fullflexible}

% --- numbers generated from the actual runs --------------------------------
% auto-generated by benchmarks/export_latex.py -- do not edit
\newcommand{\benchN}{12}
\newcommand{\benchTni}{12}
\newcommand{\benchLint}{4}
\newcommand{\benchUnit}{3}
\newcommand{\benchAsof}{3}
\newcommand{\benchMissed}{4}
\newcommand{\benchMissedList}{full-sample-zscore, winsorize-full-sample, global-feature-selection, panel-zscore}
\newcommand{\htwoN}{3}
\newcommand{\htwoDecl}{1/2/1}
\newcommand{\honeTotal}{10000}
\newcommand{\honeViol}{0}
\newcommand{\honeInterf}{4812}
\newcommand{\honeChecked}{4101}
\newcommand{\checkMedianMs}{3.0}
\newcommand{\lintBuggy}{2}
\newcommand{\lintTP}{2}
\newcommand{\lintClean}{2}
\newcommand{\lintFP}{1}


\title{\textbf{Draft \# 2 \\ Temporal Non-Interference:\\ Making Look-Ahead Bias a Compile-Time Error\\ in Quantitative Research Pipelines}}
\author{Ismail Syed\thanks{CHANGE THIS LATER NEED TO CHECK WHAT TO PUT HERE - ismailsyed2005@gmail.com}}
\date{\today}

\begin{document}
\maketitle

\begin{abstract}
Look-ahead bias---using information in a backtest that would not have been
available at decision time---is the most common and most destructive defect in
quantitative research: it produces spectacular in-sample results that evaporate
in live trading. Today it is fought with discipline (code review, point-in-time
databases, purged cross-validation); discipline does not scale and offers no
guarantee. We recast look-ahead bias as a \emph{type-safety} problem. We adapt
information-flow control---the non-interference type systems built for
security---to the \emph{time axis}: every value carries the time at which its
information first becomes \emph{available}, labels form the lattice
$(\mathbb{T},\le)$ with $\text{join}=\max$, and a pipeline is rejected unless
every decision dated $t$ depends only on information available by $t$. We define
\emph{temporal non-interference} and prove that a well-typed pipeline satisfies
it (a future input cannot change a past decision). We implement the calculus as
an embedded DSL, \tni, and evaluate it: a property-based ``perturb-the-future''
oracle finds \textbf{\honeViol{} soundness violations across \honeTotal{} random
pipelines}; a \benchN-pattern leakage benchmark on synthetic and real (equities,
crypto) data is caught \textbf{\benchTni/\benchN} while ordinary linters, unit
tests, and feature-store ASOF discipline each catch only 3--4 and miss
\benchMissed{} patterns entirely; and three realistic strategies type-check with
1--2 trusted annotations each and zero false positives. A single property
unifies look-ahead bias, full-sample preprocessing, survivorship, and improper
cross-validation as instances of the same violation.
\end{abstract}

\section{Introduction}

A quantitative strategy is a function from market history to positions. It has
look-ahead bias (a.k.a.\ temporal data leakage) when a position dated $t$ depends
on data that was not knowable until after $t$. The symptom is universal and
catastrophic: an inflated backtest that does not reproduce live. The disguises
are many---trading at the open on the same day's close; standardizing features
with full-sample statistics; centered moving averages; targets that leak into
features; universes defined by who survived to today; shuffled $K$-fold on time
series; fundamentals stamped at report-period-end rather than publication date.

The prevailing defense is \emph{vigilance}: checklists, review, point-in-time
data engineering, purged and embargoed cross-validation. This is the same posture
the systems community once took toward memory safety (``be careful with
pointers''), and it fails for the same reason: it does not scale and provides no
guarantee. Memory safety was not solved by discipline but by \emph{type
systems}. Our thesis is that look-ahead bias is likewise a type-safety problem.

\paragraph{Contributions.}
\begin{enumerate}\itemsep2pt
  \item A formal core calculus, \lt, that instantiates information-flow control
    on the time lattice, with availability-transfer rules for the operators that
    cause real leaks (\S\ref{sec:calculus}).
  \item A definition of \emph{temporal non-interference} and a soundness
    theorem: well-typed $\Rightarrow$ a future input cannot change a past
    decision (\S\ref{sec:theory}).
  \item \tni, an embedded-DSL realization with compiler-quality diagnostics,
    plus an unsound AST linter for legacy code (\S\ref{sec:system}).
  \item An evaluation establishing soundness empirically (H1), expressiveness
    (H2), and utility against baselines on real markets (H3), with quantified
    Sharpe inflation (\S\ref{sec:eval-h1h2}--\ref{sec:eval-h3}).
\end{enumerate}

\section{Background and the gap}\label{sec:background}

\paragraph{Information-flow control (IFC).}
Classic IFC assigns each value a label from a security lattice
\citep{denning1976lattice} and proves \emph{non-interference}
\citep{goguen1982security}: secret (\textsf{High}) inputs cannot influence
public (\textsf{Low}) outputs. Enforcement is a type system; the result of an
operation takes the join of its operands' labels; the guarantee is a soundness
theorem \citep{volpano1996sound}. Some downgrades are unavoidable and handled by
explicit, audited \emph{declassification} \citep{sabelfeld2003language}.

\paragraph{Point-in-time practice (finance).}
Practitioners fight leakage with point-in-time (bitemporal) databases,
feature-store ASOF joins \citep{feast}, walk-forward protocols and backtesting
hygiene \citep{arnott2019backtesting}, purged and embargoed cross-validation
\citep{lopezdeprado2018advances}, and time-ordered evaluation
\citep{bergmeir2012use}. These are methodologies and data-engineering
disciplines, not language-level static guarantees.

\paragraph{The gap.}
The two halves exist; their synthesis does not. IFC labels are always
\emph{secrecy}; temporal type systems in PL concern \emph{calendar/date
representation} \citep{bry2005catts,tc39temporal}, not information
\emph{availability}. No prior work treats time as the information-flow lattice
and look-ahead bias as a non-interference violation caught statically.
Table~\ref{tab:mapping} states the correspondence we exploit.

\begin{table}[t]\centering\small
\begin{tabular}{ll}
\toprule
Information-flow control & Temporal non-interference (this work) \\
\midrule
Label $=$ secrecy level & Label $=$ \textbf{availability time} \\
Lattice $\{\textsf{Low}\sqsubseteq\textsf{High}\}$ & Lattice $(\mathbb{T},\le)$, $\text{join}=\max$ \\
``Secret must not flow to public'' & ``\textbf{Future must not flow to a past decision}'' \\
Declassification (audited) & Publication-lag / availability assertion (audited) \\
Non-interference theorem & \textbf{Temporal non-interference theorem} \\
\bottomrule
\end{tabular}
\caption{Information-flow control, reinterpreted on the time axis.}
\label{tab:mapping}
\end{table}

\section{The \lt{} calculus}\label{sec:calculus}

\paragraph{Time and labels.}
Let $(\mathbb{T},\le)$ be a totally ordered set of timestamps, extended with
$\bot=-\infty$ (``always available,'' for constants) and $\top=+\infty$ (``never
available in-sample''). It is a lattice with $\text{join}=\max$, $\text{meet}=\min$;
\emph{later time $=$ more restricted}, mirroring \emph{more secret}. A value's
\emph{availability label} is the earliest time its information is fully
determined. For time-indexed values the label is a per-row function
$A:\mathbb{T}\to\mathbb{T}\cup\{\top\}$; a clean, strictly backward-looking
series satisfies $A(s)=s$.

\paragraph{The legality condition.}
A pipeline emits a position series $Z$. The single obligation is
\[
\forall s:\quad A_Z(s)\ \le\ s-\delta_{\text{exec}},
\]
where $\delta_{\text{exec}}\ge 0$ is an execution lag ($\delta_{\text{exec}}=1$
models ``signal at close, trade next bar'').

\paragraph{Transfer rules (excerpt).}
Each operator maps input labels to an output label.
\begin{itemize}\itemsep1pt
  \item \emph{Constant}: $\bot$. \emph{Raw observation with publication lag
    $\delta$}: $A(t)=t+\delta$.
  \item \emph{Elementwise op}: $A(s)=\bigsqcup$ of operands. \emph{Backward lag}
    $\mathrm{lag}(X,k),k\ge 0$: safe. \emph{Forward lead}: $A(s)=A_X(s+k)>s$---the
    canonical leak.
  \item \emph{Backward rolling}: join over $[s-w+1,s]$ (clean $\Rightarrow s$,
    safe). \emph{Centered/forward rolling}: window reaches $>s$---leak.
  \item \emph{Full-sample reduction} (mean/std/quantile over all time): $\top$;
    broadcasting it back and using it at $s$ leaks (full-sample normalization,
    winsorization, global feature selection).
  \item \emph{Resample labelled at period end}, used intraperiod:
    $A(s)=\text{period\_end}>s$---leak.
  \item \emph{Cross-sectional op at fixed $s$} (rank/z-score across assets): join
    over that row only---safe; a \emph{panel-pooled} op over all (asset, time)
    cells: $\top$---leak.
  \item \emph{Model fit} on data with availabilities $\le\tau$: the model has
    label $\tau$; predicting at $s$ requires $\tau\le s$. This recovers
    walk-forward / purged CV as a typing constraint; shuffled $K$-fold gives a
    model label at the sample end, leaking at every interior test point.
  \item \emph{Universe/selection}: membership from data $\le s$ is safe;
    membership referencing data $>s$ (``listed at $T_{\text{end}}$'') leaks
    (survivorship).
\end{itemize}

\paragraph{Declassification (the trusted base).}
Two primitives are the \emph{only} places trust is injected:
\texttt{observe(stream, lag=$\delta$)} (a raw source's publication lag) and
\texttt{available\_at($x,\tau$)} (assert an availability the calculus cannot
derive, e.g.\ a per-row point-in-time date). Every assertion is recorded in an
audit log; the trusted surface is small and enumerable.

\section{Temporal non-interference}\label{sec:theory}

Each input datum has a \emph{true} availability time; operators propagate both
values and true availabilities.

\begin{definition}[Temporal non-interference]
A pipeline $P$ emitting $Z$ is temporally non-interferent iff for every $s$ and
every pair of input datasets $D,D'$ that agree on all data with availability
$\le s$ (but may differ arbitrarily on data available later), $Z_D(s)=Z_{D'}(s)$.
\end{definition}

\begin{theorem}[Soundness]
If $P$ type-checks under \S\ref{sec:calculus} (so $A_Z(s)\le s$ for all $s$),
then $P$ is temporally non-interferent.
\end{theorem}

\begin{proof}[Proof sketch]
Induction on the typing derivation. Each transfer function is a monotone join, so
the inferred label of any sub-value upper-bounds the true availability of every
datum that influences it. The check $A_Z(s)\le s$ therefore guarantees only data
with true availability $\le s$ influences $Z(s)$; since $D,D'$ agree there, the
values are equal. This is the Volpano--Smith argument
\citep{volpano1996sound} with the secrecy lattice replaced by
$(\mathbb{T},\le)$.
\end{proof}

The theorem is validated empirically by the perturb-the-future test
(\S\ref{sec:eval-h1h2}), which perturbs all future inputs and checks that no past
decision moves.

\section{The \textmd{\tni} system}\label{sec:system}

\paragraph{Design: a hybrid.}
A \emph{sound core} (the embedded DSL, which carries the theorem) plus a
\emph{best-effort, explicitly unsound AST linter} for legacy code (which carries
reach). They are cleanly separated: the theorem is a statement about the core
only.

\paragraph{Build $\to$ check $\to$ run.}
The core evaluates eagerly---each operator computes values and availability
labels at construction---but every node also retains its operands
(\texttt{\_parents}) and operator name (\texttt{\_op}). We deliberately avoid a
separate lazy graph; the retained provenance recovers the one thing a graph would
buy: path-level diagnostics. The pipeline is still build $\to$ \textbf{check}
(\texttt{to\_positions} asserts $A(s)\le s-\delta_{\text{exec}}$) $\to$ run (the
backtest runs only after the check passes). \emph{Limitation:} labels are
computed by executing the builder, so data-dependent control flow is not
modelled; quant pipelines are overwhelmingly straight-line dataflow, and the
perturb oracle is the independent check that the labels are sound regardless.

\paragraph{Diagnostics.}
On a violation, \tni{} reports the offending operator, decision time, inferred vs.\
required label, and the dependency path, descending only through illegal operands
and marking where future information enters:

\begin{lstlisting}[style=plain]
look-ahead bias: position at t=0 depends on information
  available only at A=2 (required A <= 0)
  dependency path (future -> decision):
    sign: A(0)=2 (need <= 0)
        roll_back3.mean: A(0)=2 (need <= 0)
            shift(k=-2): A(0)=2 (need <= 0) <== future info enters here
\end{lstlisting}

Type-checking is cheap: the median check time over the benchmark pipelines is
\checkMedianMs~ms, so it is a build-time gate, not a runtime cost.

\paragraph{Unsound AST linter.}
A syntactic pass over ordinary pandas/numpy/sklearn source flags red flags
(negative \texttt{.shift()}, \texttt{rolling(center=True)}, a shuffled
\texttt{KFold} or \texttt{train\_test\_split}, a scaler \texttt{fit\_transform},
full-sample statistics in arithmetic). It is a bug-finder with false positives
and negatives by design (\S\ref{sec:eval-h3}).

\section{Evaluation: soundness (H1) and expressiveness (H2)}\label{sec:eval-h1h2}

\paragraph{H1 --- soundness, empirically.}
A generator builds random pipelines from safe and unsafe operators; the oracle
perturbs every input with true availability $>s^\star$ and checks that no
decision at $s\le s^\star$ changes. Over \textbf{\honeTotal{} random pipelines}:
\textbf{\honeViol{} violations} (no pipeline both type-checks and depends on the
future). The test is non-vacuous---\honeInterf{} generated pipelines genuinely
interfere and the checker rejects every one. Two transfer-function bugs were
found and fixed this way during development (rolling and cross-sectional joins
must be taken over \emph{finite} contributors only).

\paragraph{H2 --- expressiveness.}
Three realistic, correct strategies, end-to-end on real data
(Table~\ref{tab:h2}). All type-check (\textbf{zero false positives}) with 1--2
trusted annotations each (declassifications \htwoDecl); every other label is
inferred.

\begin{table}[t]
\centering\small
\begin{tabular}{llcccr}
\toprule
Strategy & universe & checks & decl. & ops & Sharpe \\
\midrule
cross-sectional momentum & crypto x8 & \cmark & 1 & 6 & $+0.33$ \\
PIT long-horizon reversal & equities x20 & \cmark & 2 & 8 & $-0.16$ \\
walk-forward GBM (embargoed CV) & crypto: BTC/USD & \cmark & 1 & 2 & $+1.00$ \\
\bottomrule
\end{tabular}
\caption{H2 case studies: all type-check (zero false positives) with 1--2 trusted declassifications each.}
\label{tab:h2}
\end{table}

\section{Evaluation: utility (H3) and applicability}\label{sec:eval-h3}

\paragraph{The benchmark.}
\benchN{} leakage patterns, each as a buggy/correct pair, run on
synthetic-with-structure data (Table~\ref{tab:bench-synth}) and real equities
(Table~\ref{tab:bench-eq}) and crypto (Table~\ref{tab:bench-crypto}). We report
whether \tni{} rejects the buggy pipeline and accepts the correct one, and the
in-sample Sharpe the buggy pipeline \emph{would} have reported had the leak
slipped past review. Baselines: a syntactic \emph{linter}, a \emph{unit-test}
Sharpe-sanity check ($|\text{Sharpe}|>5$), and \emph{ASOF}/point-in-time
discipline.

\begin{table}[t]
\centering\small
\begin{tabular}{lccccrr}
\toprule
Leakage pattern & \textsc{tni} & lint & unit & ASOF & Sharpe$^\dagger$ & honest \\
\midrule
close-to-open & \cmark & \xmark & \cmark & \xmark & $+22.03$ & $+0.65$ \\
target-leakage & \cmark & \cmark & \cmark & \xmark & $+22.03$ & $+0.65$ \\
full-sample-zscore & \cmark & \xmark & \xmark & \xmark & $+0.16$ & $+0.08$ \\
centered-rolling & \cmark & \cmark & \xmark & \xmark & $+4.58$ & $+0.26$ \\
winsorize-full-sample & \cmark & \xmark & \xmark & \xmark & $+0.65$ & $+0.65$ \\
shuffled-kfold-cv & \cmark & \cmark & \xmark & \xmark & $+0.86$ & $+0.93$ \\
resampling-lookahead & \cmark & \xmark & \xmark & \cmark & $+4.52$ & $+0.79$ \\
global-feature-selection & \cmark & \xmark & \xmark & \xmark & $+0.30$ & $+0.25$ \\
publication-lag & \cmark & \xmark & \xmark & \cmark & $+0.32$ & $+0.57$ \\
panel-zscore & \cmark & \xmark & \xmark & \xmark & $+0.19$ & $+0.15$ \\
survivorship-universe & \cmark & \xmark & \xmark & \cmark & $+1.40$ & $+0.24$ \\
target-leakage-panel & \cmark & \cmark & \cmark & \xmark & $+119.28$ & $+0.15$ \\
\bottomrule
\end{tabular}
\caption{Leakage benchmark on seeded synthetic data with genuine predictable structure. $\dagger$: in-sample Sharpe the buggy pipeline would report if undetected.}
\label{tab:bench-synth}
\end{table}

\tni{} catches \textbf{\benchTni/\benchN}; the linter \benchLint/\benchN, the
unit-test \benchUnit/\benchN, ASOF \benchAsof/\benchN. The same \benchTni/\benchN{}
holds on equities and crypto. Crucially the baselines catch \emph{different}
small subsets, so even their union misses \benchMissed{} patterns
---\benchMissedList---that \tni{} rejects by construction. This is the rebuttal
to ``isn't this just ASOF joins?'': ASOF discipline addresses timestamp
alignment (publication-lag, survivorship, resampling) but is blind to full-sample
statistical leakage and CV mistakes.

\begin{table}[t]
\centering\small
\begin{tabular}{lccccrr}
\toprule
Leakage pattern & \textsc{tni} & lint & unit & ASOF & Sharpe$^\dagger$ & honest \\
\midrule
close-to-open & \cmark & \xmark & \cmark & \xmark & $+15.84$ & $+0.43$ \\
target-leakage & \cmark & \cmark & \cmark & \xmark & $+15.84$ & $+0.43$ \\
full-sample-zscore & \cmark & \xmark & \xmark & \xmark & $+0.19$ & $+0.36$ \\
centered-rolling & \cmark & \cmark & \xmark & \xmark & $+3.80$ & $+0.24$ \\
winsorize-full-sample & \cmark & \xmark & \xmark & \xmark & $+0.43$ & $+0.42$ \\
shuffled-kfold-cv & \cmark & \cmark & \xmark & \xmark & $+0.95$ & $+0.90$ \\
resampling-lookahead & \cmark & \xmark & \xmark & \cmark & $+4.49$ & $+0.12$ \\
global-feature-selection & \cmark & \xmark & \xmark & \xmark & $-0.07$ & $+0.06$ \\
publication-lag & \cmark & \xmark & \xmark & \cmark & $+0.36$ & $+0.16$ \\
panel-zscore & \cmark & \xmark & \xmark & \xmark & $-0.07$ & $+0.10$ \\
survivorship-universe & \cmark & \xmark & \xmark & \cmark & $+1.04$ & $+0.83$ \\
target-leakage-panel & \cmark & \cmark & \cmark & \xmark & $+32.37$ & $+0.10$ \\
\bottomrule
\end{tabular}
\caption{Leakage benchmark on US equities (20 names, 2010--2023).}
\label{tab:bench-eq}
\end{table}
\begin{table}[t]
\centering\small
\begin{tabular}{lccccrr}
\toprule
Leakage pattern & \textsc{tni} & lint & unit & ASOF & Sharpe$^\dagger$ & honest \\
\midrule
close-to-open & \cmark & \xmark & \cmark & \xmark & $+15.89$ & $-0.39$ \\
target-leakage & \cmark & \cmark & \cmark & \xmark & $+15.89$ & $-0.39$ \\
full-sample-zscore & \cmark & \xmark & \xmark & \xmark & $-0.22$ & $-0.01$ \\
centered-rolling & \cmark & \cmark & \xmark & \xmark & $+3.22$ & $+0.24$ \\
winsorize-full-sample & \cmark & \xmark & \xmark & \xmark & $-0.39$ & $-0.40$ \\
shuffled-kfold-cv & \cmark & \cmark & \xmark & \xmark & $+0.02$ & $-0.39$ \\
resampling-lookahead & \cmark & \xmark & \xmark & \cmark & $+4.23$ & $+0.12$ \\
global-feature-selection & \cmark & \xmark & \xmark & \xmark & $+0.12$ & $-0.62$ \\
publication-lag & \cmark & \xmark & \xmark & \cmark & $+0.17$ & $+0.06$ \\
panel-zscore & \cmark & \xmark & \xmark & \xmark & $+0.48$ & $+0.41$ \\
survivorship-universe & \cmark & \xmark & \xmark & \cmark & $+0.41$ & $+0.15$ \\
target-leakage-panel & \cmark & \cmark & \cmark & \xmark & $+20.28$ & $+0.41$ \\
\bottomrule
\end{tabular}
\caption{Leakage benchmark on crypto (8 pairs, daily).}
\label{tab:bench-crypto}
\end{table}

Two honest observations. (i) Not every leak inflates Sharpe on every dataset---
full-sample standardization and winsorization barely move a sign-based signal's
Sharpe; their harm is more in risk estimates. \tni{} rejects them regardless of
whether they happen to pay off on a given sample. (ii) The dramatic numbers
(close-to-open, target leakage) come from leaks that mechanically capture the
realized return.

\paragraph{Applicability (AST linter).}
Over a curated corpus of representative scripts, the linter flags
\textbf{\lintTP/\lintBuggy} buggy scripts and produces \lintFP{} false positive
on a clean script: it flags a forward-shifted \emph{target} ($y=$ next return),
which is legitimate---a purely syntactic tool cannot distinguish a leaked feature
from a label. The sound core has no such false positive (cf.\ H2: \htwoN/\htwoN).
The linter's misses and false alarms are precisely the argument for why the core,
not the linter, carries the guarantee.

\section{Related work}\label{sec:related}

IFC and non-interference type systems---the lattice model
\citep{denning1976lattice}, non-interference \citep{goguen1982security}, sound
flow typing \citep{volpano1996sound}, the language-based survey
\citep{sabelfeld2003language}, and dynamic enforcement via faceted values
\citep{austin2012multiple} and secure multi-execution
\citep{devriese2010noninterference}---all use secrecy labels, not time.
Point-in-time / bitemporal data and feature stores \citep{feast}---data
engineering, no language guarantee. Purged/embargoed CV
\citep{lopezdeprado2018advances} and time-ordered evaluation
\citep{bergmeir2012use}---methodology, recovered here as a typing constraint;
survivorship bias has long been documented empirically
\citep{brown1992survivorship}. Backtesting protocols
\citep{arnott2019backtesting} prescribe discipline rather than enforce it.
Temporal types \citep{bry2005catts,tc39temporal}---calendar representation, not
availability.

\section{Limitations and future work}\label{sec:limitations}

\begin{itemize}\itemsep1pt
  \item \emph{Mechanized proof.} The soundness proof is on paper plus the perturb
    oracle; a Coq/Lean formalization is future work.
  \item \emph{Control flow.} Eager evaluation does not model data-dependent
    branching (host-language \texttt{if} on values escapes into an implicit
    flow); a typed conditional that joins the guard's availability, or a static
    front-end over a Python subset, would extend coverage soundly.
  \item \emph{Trusted base.} \texttt{observe}/\texttt{available\_at} are believed,
    not checked---publication-lag is caught only once the true lag is declared.
    This is by design (the audited base) but must be stated.
  \item \emph{Coverage.} The operator set covers the benchmark and case studies,
    not all of pandas; the AST linter is unsound.
  \item Intraday/event-time semantics, restatements/bitemporality, multi-source
    clock skew, and IDE integration are open.
\end{itemize}

\section{Conclusion}\label{sec:conclusion}

Treating availability time as an information-flow lattice turns look-ahead bias
from a discipline problem into a type error, caught before any backtest runs, with
a soundness guarantee. One property---temporal non-interference---unifies
look-ahead bias, full-sample preprocessing, survivorship, and improper
cross-validation. The prototype catches a \benchN-pattern benchmark in full on
real markets where conventional tooling catches a third, while accepting realistic
strategies with negligible annotation burden. Leakage, like memory safety, is
better prevented by construction than by vigilance.

\bibliographystyle{plainnat}
\bibliography{references}

\end{document}
